(1) $\lambda=2\pi\int_0^a\rho_0(1-r^2/a^2)r\,dr=\pi\rho_0a^2/2$。Gauss定律给出 $E_r(r<a)=\rho_0[r/2-r^3/(4a^2)]/\varepsilon$，$E_r(a<r<b)=\rho_0a^2/(4\varepsilon r)$；金属筒屏蔽外场。
(2) 由 $\varphi(b)=0$ 得 $\varphi(a)=\rho_0a^2\ln(b/a)/(4\varepsilon)$。柱内 $\varphi(r)=\varphi(a)+\rho_0[(a^2-r^2)/4-(a^4-r^4)/(16a^2)]/\varepsilon$，特别地 $\varphi(0)=\varphi(a)+3\rho_0a^2/(16\varepsilon)$，并满足 $-\varphi'=E_r$。
(3) 单位长度电场能 $U'=\pi\varepsilon\int_0^b E_r^2r\,dr$。柱内 $U'_{\rm in}=11\pi\rho_0^2a^4/(384\varepsilon)$，柱外 $U'_{\rm out}=\pi\rho_0^2a^4\ln(b/a)/(16\varepsilon)$；两项均非负，单位为焦耳每米。
(4) 中心 $E_r\sim \rho_0r/(2\varepsilon)$ 无奇点；$E_r(a^-)=E_r(a^+)=\rho_0a/(4\varepsilon)$，符合无面自由电荷边界条件。$b\to a^+$ 时外侧储能趋零，但内部空间电荷储能仍非零。
